% !TEX root = ../../../main.tex

\subsection{视频读取与帧级检测结果缓存}

句子视频识别的输入来自移动端提交的本地视频文件。
在 HearBridge 系统中，用户可以在移动端选择一段待识别视频，并通过句子视频识别链路将视频上传到服务端。
Python 手势识别服务接收上传文件后，会先完成文件校验、临时保存和视频读取等前置处理。
这些步骤为后续关键点检测、滑动窗口预测和识别结果生成提供基础。

\WhuAutoFigure{0.82\textwidth}{0.42\textheight}{figures/chapter06/frame-cache-flow.png}{句子视频识别帧级检测缓存流程图}{fig:chapter06-frame-cache-flow}

如图~\ref{fig:chapter06-frame-cache-flow} 所示，视频读取与帧级检测结果缓存是句子视频识别流程的前置环节。
该环节的目标不是直接完成词汇分类，而是将上传视频转换为后续识别流程可以复用的帧序列和 MediaPipe 检测结果序列。
系统在这一阶段完成视频文件保存、视频元信息读取、逐帧解码、关键点检测以及检测状态统计等工作。

服务端接收到上传视频后，首先会对上传文件进行基本校验。
校验内容包括文件名后缀、文件是否为空以及文件大小是否超过限制等。
当前句子视频识别链路允许处理 mp4、mov、avi、mkv 等常见视频格式。
如果文件类型不符合要求，或者上传内容为空，系统会直接返回错误信息，避免无效输入继续进入后续识别流程。
通过前置校验，可以减少异常视频文件对服务端推理流程的影响。

校验通过后，系统会为本次识别请求创建独立的临时目录。
临时目录通常按照本次请求的 request\_id 进行区分。
随后，系统将上传视频保存为该目录下的输入文件。
这种做法可以避免不同识别请求之间的文件相互覆盖，也便于在一次识别流程结束后清理本次请求产生的临时文件。
在调试场景下，系统也可以保留临时目录，便于开发者查看上传视频、帧级中间结果和识别输出文件。

视频保存完成后，系统使用 OpenCV 打开待识别视频。
如果视频无法被正常打开，说明上传文件可能损坏、格式不受支持或内容不完整。
此时系统会终止识别流程并返回错误。
当视频能够正常打开时，系统会读取视频的基础元信息。
这些元信息包括帧宽、帧高、帧率和总帧数等。
视频元信息不会直接作为模型输入参与分类，但会被写入最终识别结果，用于辅助分析滑动窗口位置和识别结果是否合理。

完成元信息读取后，系统通过 VideoCapture 顺序读取视频中的全部帧。
每一帧都会被保存到内存中的 frames 列表中。
如果读取结束后帧列表为空，系统同样会终止识别流程。
这是因为后续滑动窗口构造需要基于完整的视频帧序列进行。
只有当系统成功获得有效帧序列后，才会进入帧级关键点检测阶段。

在获得视频帧序列后，系统使用 MediaPipe Holistic 对每一帧进行人体姿态和左右手关键点检测。
Holistic 检测结果中包含 pose\_landmarks、left\_hand\_landmarks 和 right\_hand\_landmarks。
其中，pose\_landmarks 用于描述人体姿态关键点，left\_hand\_landmarks 和 right\_hand\_landmarks 用于描述左右手关键点。
这些检测结果是后续构造 $20\times232$ plus 特征的原始点位来源。

系统不会在每一个滑动窗口中重复执行 MediaPipe 检测。
相反，系统会先对整段视频的每一帧执行一次 Holistic 检测。
每一帧的检测结果会按照帧顺序加入 results 列表。
results 列表与 frames 列表一一对应。
这样，后续任意滑动窗口只需要根据帧下标从 results 中读取缓存结果，而不需要重新对同一帧运行 MediaPipe。
这种缓存方式可以显著减少重复计算，提高句子视频识别阶段的推理效率。

在缓存 MediaPipe 检测结果的同时，系统还会记录每一帧的检测状态。
pose\_flags 用于表示当前帧是否检测到人体姿态关键点。
hand\_flags 用于表示当前帧是否至少检测到一只手。
当左手或右手任意一侧关键点存在时，该帧的 hand\_flags 即为 True。
这两个标志不直接作为分类标签使用，但能够反映输入视频的关键点检测质量。

完成整段视频的帧级检测后，系统会根据 pose\_flags 和 hand\_flags 计算检测覆盖率。
pose\_ratio 表示整段视频中成功检测到人体姿态关键点的帧比例。
any\_hand\_ratio 表示整段视频中至少检测到一只手的帧比例。
这两个指标会被写入识别结果中。
如果 pose\_ratio 或 any\_hand\_ratio 较低，说明当前视频中人体姿态或手部关键点检测不稳定。
在这种情况下，后续滑动窗口特征构造和词级预测结果也可能受到影响。

帧级检测缓存还为后续动态差分特征构造提供支持。
在滑动窗口阶段，系统会从 results 中取出窗口范围内的检测结果。
对于窗口内的每一个采样帧，系统根据缓存的 pose 和 hand 关键点构造单帧 plus 特征。
同时，系统会在窗口内部维护上一帧的动态源向量，用于计算相邻帧之间的 dynamic\_delta 特征。
窗口第一帧没有上一帧可供比较，因此其动态差分可以视为从零开始。
这种处理方式与训练阶段对独立词级样本构造 20 帧特征的方式保持一致。

通过上述设计，系统将视频读取、帧级关键点检测和滑动窗口特征构造进行了解耦。
视频读取阶段负责获得帧序列和元信息。
帧级检测阶段负责一次性缓存每一帧的 MediaPipe 结果。
滑动窗口阶段则复用这些缓存结果构造多个候选窗口的模型输入。
这种流程既减少了重复检测开销，也使后续识别结果能够结合视频元信息、检测覆盖率和窗口位置进行分析。
最终，帧序列、检测缓存、pose\_ratio、any\_hand\_ratio 等信息会共同支撑后续的窗口预测、词段合并和句子级识别结果生成。